% !TEX program = xelatex
% 市场研究报告模板
% 适用于50页以上综合市场报告的专业排版
% 配合 market_research.sty 样式包使用

\documentclass[11pt,letterpaper]{report}
\usepackage{market_research}

% ============================================================================
% 文档元数据 - 请根据需要自定义
% ============================================================================
\newcommand{\reporttitle}{[MARKET NAME]}
\newcommand{\reportsubtitle}{综合市场分析报告}
\newcommand{\reportdate}{\today}
\newcommand{\reportauthor}{市场情报部}
\newcommand{\reportclassification}{机密}

% ============================================================================
% PDF 元数据
% ============================================================================
\hypersetup{
    pdftitle={\reporttitle{} - \reportsubtitle{}},
    pdfauthor={\reportauthor{}},
    pdfsubject={市场研究报告},
    pdfkeywords={市场研究, 市场分析, 竞争格局, 战略分析}
}

% ============================================================================
% 文档开始
% ============================================================================
\begin{document}

% ============================================================================
% 封面页
% ============================================================================
% 如需使用封面图片，请将空括号替换为图片路径：
% \makemarketreporttitle{\reporttitle}{\reportsubtitle}{../figures/cover_image.png}{\reportdate}{\reportauthor}

\begin{titlepage}
\centering
\vspace*{2cm}

{\Huge\bfseries\color{primaryblue} \reporttitle\\[0.5cm]}
{\LARGE\bfseries \reportsubtitle\\[1.5cm]}

% 视觉：生成封面/主图
% python skills/generate-image/scripts/generate_image.py "专业执行摘要信息图，现代数据可视化风格展示关键指标，蓝绿配色方案，简洁极简设计" --output figures/cover_image.png
% 图片生成后取消下方注释：
% \includegraphics[width=\textwidth]{../figures/cover_image.png}\\[1.5cm]

\vspace{4cm}

{\Large\bfseries 综合市场研究报告\\[0.5cm]}
{\large 面向商业决策的战略情报\\[3cm]}

{\large
\textbf{日期：} \reportdate\\[0.3cm]
\textbf{编制：} \reportauthor\\[0.3cm]
\textbf{密级：} \reportclassification\\[0.3cm]
\textbf{报告类型：} 完整市场分析
}

\vfill

{\footnotesize
\textit{本报告包含基于公开数据和自有研究的市场情报与战略分析。所有来源均已标注并可独立验证。}
}

\end{titlepage}

% ============================================================================
% 前置内容
% ============================================================================
\pagenumbering{roman}

% 目录
\tableofcontents
\newpage

% 插图列表
\listoffigures
\newpage

% 表格列表
\listoftables
\newpage

% ============================================================================
% 正文内容
% ============================================================================
\pagenumbering{arabic}

% ============================================================================
% 执行摘要（2-3页）
% ============================================================================
\chapter{执行摘要}

\section{报告概述}

本综合市场分析考察了 \reporttitle{} 市场，为投资者、高管和战略规划者提供战略情报。本报告综合了来自权威来源的数据，包括市场研究机构、监管机构、行业协会和企业调查。

% 视觉：执行摘要信息图
% python skills/scientific-schematics/scripts/generate_schematic.py "执行摘要信息图展示4个关键指标：市场规模$XX.XB(2024)、复合年均增长率XX.X%、前3大参与者、关键趋势。蓝底白字，专业布局" -o figures/exec_summary_infographic.png --doc-type report

\subsection{市场快照}

\begin{marketdatabox}[市场快照：\reporttitle{}]
\begin{itemize}[leftmargin=*]
    \item \textbf{当前市场规模（2024年）：} \marketsize{X.XX十亿美元}
    \item \textbf{预测市场规模（2034年）：} \marketsize{XX.XX十亿美元}
    \item \textbf{复合年增长率（CAGR）：} \growthrate{XX.X}
    \item \textbf{增长倍数：} 10年内增长X倍
    \item \textbf{最大细分市场：} [细分市场名称]（XX\% 市场份额）
    \item \textbf{增长最快的地区：} [地区]（\growthrate{XX.X} CAGR）
    \item \textbf{当前企业采用率：} XX\%
\end{itemize}
\end{marketdatabox}

\subsection{投资论点}

多重市场催化剂的汇聚为 \reporttitle{} 市场的投资和战略行动创造了引人注目的机会：

\begin{keyinsightbox}[关键投资驱动因素]
\begin{enumerate}
    \item \textbf{[驱动因素1：]} [简要说明该驱动因素如何创造机会]
    \item \textbf{[驱动因素2：]} [简要说明]
    \item \textbf{[驱动因素3：]} [简要说明]
    \item \textbf{[驱动因素4：]} [简要说明]
\end{enumerate}
\end{keyinsightbox}

\subsection{主要发现}

\paragraph{市场动态}
[总结关于市场规模、增长和动态的最重要发现。包含3-5个关键统计数据及其来源。]

\paragraph{竞争格局}
[总结竞争动态、市场集中度和关键参与者。包含主要参与者的市场份额。]

\paragraph{增长驱动因素}
[总结推动市场增长的主要因素及其预期影响。]

\paragraph{风险因素]
[总结可能影响市场发展的关键风险。]

\subsection{战略建议}

基于本报告中的综合分析，我们建议采取以下战略行动：

\begin{recommendationbox}[顶级战略建议]
\begin{enumerate}
    \item \textbf{[建议1：]} [面向行动的建议及预期成果]
    \item \textbf{[建议2：]} [面向行动的建议及预期成果]
    \item \textbf{[建议3：]} [面向行动的建议及预期成果]
    \item \textbf{[建议4：]} [面向行动的建议及预期成果]
    \item \textbf{[建议5：]} [面向行动的建议及预期成果]
\end{enumerate}
\end{recommendationbox}

% ============================================================================
% 第1章：市场概述与定义（4-5页）
% ============================================================================
\chapter{市场概述 \& 定义}

\section{市场定义}

[对所分析市场提供清晰、全面的定义，包括：
- 涵盖哪些产品/服务
- 明确排除哪些内容
- 该市场与相邻市场的关系
- 适用的行业分类代码（NAICS、SIC）]

\begin{calloutbox}[市场定义]
\reporttitle{} 市场涵盖[全面定义]。这包括[包含的要素]并排除[排除的要素]。
\end{calloutbox}

\subsection{范围与边界}

\paragraph{地理范围}
[定义分析的地理边界——全球、区域或特定国家。]

\paragraph{产品/服务范围}
[定义市场定义中包含的产品和服务。]

\paragraph{时间范围}
[说明分析的历史时期和预测期间。]

\subsection{市场分类}

[提供详细的市场分类和分类体系，包括：
- 市场细分
- 子细分
- 类别]

\section{产业生态系统}

% 视觉：产业生态系统图
% python skills/scientific-schematics/scripts/generate_schematic.py "产业生态价值链图，从左到右：原材料/投入→制造/加工→分销→终端用户。各阶段标注关键参与者。蓝色方框箭头连接" -o figures/industry_ecosystem.png --doc-type report

[描述产业生态系统和价值链，包括：
- 关键利益相关者及其角色
- 利益相关者之间的关系
- 各阶段的价值创造
- 信息和资金流]

\begin{figure}[htbp]
\centering
% \includegraphics[width=0.95\textwidth]{../figures/industry_ecosystem.png}
\caption{产业生态系统与价值链}
\label{fig:ecosystem}
\end{figure}

\subsection{关键利益相关者}

[描述各类利益相关者：]

\paragraph{供应商/厂商}
[上游供应商及其角色的描述。]

\paragraph{制造商/服务商}
[核心市场参与者的描述。]

\paragraph{分销商/渠道}
[分销和市场进入渠道的描述。]

\paragraph{终端用户/客户}
[客户细分及其需求的描述。]

\paragraph{监管机构与行业组织}
[监管环境和行业协会的描述。]

\section{市场结构}

% 视觉：市场结构图
% python skills/scientific-schematics/scripts/generate_schematic.py "市场结构图展示产业层次：中心为核心市场，周围为邻近市场，底层为使能技术，顶层为监管框架。同心矩形" -o figures/market_structure.png --doc-type report

[描述市场结构：]

\subsection{市场集中度}

[使用以下指标分析市场集中度：
- 赫芬达尔-赫希曼指数（HHI）
- CR4/CR8集中度
- 市场碎片化评估]

\subsection{产业生命周期阶段}

[识别市场所处的生命周期阶段：
- 导入期
- 成长期
- 成熟期
- 衰退期]

\section{历史背景}

[提供市场的历史背景：
- 市场何时出现？
- 市场发展的关键里程碑
- 重大行业变革和颠覆
- 市场随时间如何演变？]

% ============================================================================
% 第2章：市场规模与增长分析（6-8页）
% ============================================================================
\chapter{市场规模 \& 增长分析}

\section{总可寻址市场（TAM）}

总可寻址市场代表在获得100\%市场份额的情况下可获得的总收入机会。基于多个研究来源的综合分析：

% 视觉：市场增长轨迹
% python skills/scientific-schematics/scripts/generate_schematic.py "柱状图展示2020至2034年市场增长。历史柱形2020-2024深蓝色，预测柱形2025-2034浅蓝色。Y轴十亿美元，X轴年份。标注复合年均增长率。标题：[MARKET]市场增长趋势" -o figures/market_growth_trajectory.png --doc-type report

\begin{table}[htbp]
\centering
\caption{全球 \reporttitle{} 市场预测（2024-2034）}
\begin{tabular}{@{}lrrr@{}}
\toprule
\textbf{年份} & \textbf{市场规模（美元）} & \textbf{同比增长} & \textbf{CAGR} \\
\midrule
2024 & \$X.XX十亿 & -- & -- \\
\rowcolor{tablealt} 2025 & \$X.XX十亿 & XX.X\% & XX.X\% \\
2026 & \$X.XX十亿 & XX.X\% & XX.X\% \\
\rowcolor{tablealt} 2027 & \$X.XX十亿 & XX.X\% & XX.X\% \\
2028 & \$X.XX十亿 & XX.X\% & XX.X\% \\
\rowcolor{tablealt} 2029 & \$X.XX十亿 & XX.X\% & XX.X\% \\
2030 & \$X.XX十亿 & XX.X\% & XX.X\% \\
\rowcolor{tablealt} 2031 & \$X.XX十亿 & XX.X\% & XX.X\% \\
2032 & \$X.XX十亿 & XX.X\% & XX.X\% \\
\rowcolor{tablealt} 2033 & \$X.XX十亿 & XX.X\% & XX.X\% \\
2034 & \$X.XX十亿 & XX.X\% & XX.X\% \\
\bottomrule
\end{tabular}
\label{tab:tam_projections}
\end{table}

\textbf{来源：} [主要研究来源，年份]

\begin{figure}[htbp]
\centering
% \includegraphics[width=0.9\textwidth]{../figures/market_growth_trajectory.png}
\caption{市场增长轨迹（2020-2034）}
\label{fig:market_growth}
\end{figure}

\subsection{历史增长分析}

[分析过去5-10年的历史市场表现：
- 历史CAGR
- 关键增长期及驱动因素
- 重大事件（衰退、颠覆等）的影响
- 与整体经济增长的比较]

\subsection{增长预测方法论}

[解释增长预测背后的方法论：
- 关键假设
- 数据来源
- 建模方法
- 置信区间]

\section{可服务可寻址市场（SAM）}

可服务可寻址市场代表在当前产品供应、地理覆盖和监管约束下可服务的TAM部分。

% 视觉：TAM/SAM/SOM图
% python skills/scientific-schematics/scripts/generate_schematic.py "TAM SAM SOM同心圆图。外圈TAM $XXB全部可及市场，中圈SAM $XXB可服务可及市场，内圈SOM $XXB可获得可服务市场。箭头指向各圈标注。专业蓝色配色" -o figures/tam_sam_som.png --doc-type report

\begin{figure}[htbp]
\centering
% \includegraphics[width=0.8\textwidth]{../figures/tam_sam_som.png}
\caption{TAM/SAM/SOM 市场机会分解}
\label{fig:tam_sam_som}
\end{figure}

\begin{table}[htbp]
\centering
\caption{SAM内的市场细分（2024年 vs 2034年）}
\begin{tabular}{@{}lrrrr@{}}
\toprule
\textbf{细分市场} & \textbf{2024年价值} & \textbf{2034年价值} & \textbf{CAGR} & \textbf{份额} \\
\midrule
细分市场A & \$X.XX十亿 & \$XX.XX十亿 & XX.X\% & XX\% \\
\rowcolor{tablealt} 细分市场B & \$X.XX十亿 & \$XX.XX十亿 & XX.X\% & XX\% \\
细分市场C & \$X.XX十亿 & \$XX.XX十亿 & XX.X\% & XX\% \\
\rowcolor{tablealt} 细分市场D & \$X.XX十亿 & \$XX.XX十亿 & XX.X\% & XX\% \\
细分市场E & \$X.XX十亿 & \$XX.XX十亿 & XX.X\% & XX\% \\
\midrule
\textbf{SAM总计} & \textbf{\$X.XX十亿} & \textbf{\$XX.XX十亿} & \textbf{XX.X\%} & \textbf{100\%} \\
\bottomrule
\end{tabular}
\label{tab:sam_segments}
\end{table}

\section{可服务可获取市场（SOM）}

可服务可获取市场代表基于竞争动态、市场进入能力和战略定位的现实市场份额获取。

\begin{keyinsightbox}[SOM预测（2034年）]
\textbf{保守情景（XX\% 市场份额）：} \marketsize{X.X十亿美元}
\begin{itemize}
    \item 假设竞争性市场中存在多个主要参与者
    \item 典型的XX-XX个月企业销售周期
    \item 聚焦于[特定细分市场]
\end{itemize}

\textbf{基准情景（XX\% 市场份额）：} \marketsize{X.X十亿美元}
\begin{itemize}
    \item 在关键细分市场获取先发优势
    \item 强劲的产品-市场匹配
    \item 建立的合作生态系统
\end{itemize}

\textbf{乐观情景（XX\% 市场份额）：} \marketsize{X.X十亿美元}
\begin{itemize}
    \item 市场领导地位
    \item 平台效应和网络优势
    \item 专有优势和护城河
\end{itemize}
\end{keyinsightbox}

\section{区域市场分析}

% 视觉：区域分布
% python skills/scientific-schematics/scripts/generate_schematic.py "饼图或树状图展示区域市场分布。北美XX%，欧洲XX%，亚太XX%，拉丁美洲XX%，中东非洲XX%。各区域不同颜色。含百分比和金额" -o figures/regional_breakdown.png --doc-type report

\begin{figure}[htbp]
\centering
% \includegraphics[width=0.8\textwidth]{../figures/regional_breakdown.png}
\caption{各区域市场规模（2024年）}
\label{fig:regional_breakdown}
\end{figure}

\begin{table}[htbp]
\centering
\caption{区域市场规模与增长}
\begin{tabular}{@{}lrrrr@{}}
\toprule
\textbf{地区} & \textbf{2024年规模} & \textbf{份额} & \textbf{CAGR} & \textbf{2034年规模} \\
\midrule
北美 & \$X.XX十亿 & XX.X\% & XX.X\% & \$XX.XX十亿 \\
\rowcolor{tablealt} 欧洲 & \$X.XX十亿 & XX.X\% & XX.X\% & \$XX.XX十亿 \\
亚太 & \$X.XX十亿 & XX.X\% & XX.X\% & \$XX.XX十亿 \\
\rowcolor{tablealt} 拉丁美洲 & \$X.XX十亿 & XX.X\% & XX.X\% & \$XX.XX十亿 \\
中东 \& 非洲 & \$X.XX十亿 & XX.X\% & XX.X\% & \$XX.XX十亿 \\
\midrule
\textbf{全球总计} & \textbf{\$X.XX十亿} & \textbf{100\%} & \textbf{XX.X\%} & \textbf{\$XX.XX十亿} \\
\bottomrule
\end{tabular}
\label{tab:regional_market}
\end{table}

\subsection{北美}
[北美市场的详细分析，包括美国和加拿大的具体情况。]

\subsection{欧洲}
[欧洲市场的详细分析，包括主要国家的细分。]

\subsection{亚太}
[亚太市场的详细分析，聚焦中国、日本、印度和新兴市场。]

\subsection{其他地区}
[拉丁美洲、中东和非洲市场的分析。]

\section{细分市场分析}

% 视觉：细分市场增长比较
% python skills/scientific-schematics/scripts/generate_schematic.py "横向柱状图对比细分领域增长率。Y轴为细分领域，X轴为复合年均增长率。柱形颜色从绿(最高)到蓝(最低)。每个柱形带数据标签" -o figures/segment_growth_comparison.png --doc-type report

\begin{figure}[htbp]
\centering
% \includegraphics[width=0.9\textwidth]{../figures/segment_growth.png}
\caption{细分市场增长率比较（CAGR 2024-2034）}
\label{fig:segment_growth}
\end{figure}

[提供每个市场细分的详细分析，包括：
- 当前规模和市场份额
- 增长轨迹
- 各细分市场的关键驱动因素
- 细分市场内的竞争动态]

% ============================================================================
% 第3章：行业驱动因素与趋势（5-6页）
% ============================================================================
\chapter{行业驱动因素 \& 趋势}

\section{主要增长驱动因素}

[识别和分析推动市场增长的关键因素：]

% 视觉：驱动因素影响矩阵
% python skills/scientific-schematics/scripts/generate_schematic.py "2x2矩阵展示市场驱动因素。X轴影响低到高，Y轴可能性低到高。绘制8-10个驱动因素圆点。右上'关键驱动'，右下'密切关注'，左上'监控'，左下'次要优先'。专业蓝绿配色" -o figures/market_drivers_matrix.png --doc-type report

\begin{figure}[htbp]
\centering
% \includegraphics[width=0.85\textwidth]{../figures/driver_impact_matrix.png}
\caption{市场驱动因素影响评估矩阵}
\label{fig:driver_matrix}
\end{figure}

\subsection{驱动因素1：[名称]}
[对该驱动因素的详细分析，包括：
- 如何影响市场
- 量化影响
- 影响时间线
- 支持证据和数据]

\subsection{驱动因素2：[名称]}
[详细分析]

\subsection{驱动因素3：[名称]}
[详细分析]

\subsection{驱动因素4：[名称]}
[详细分析]

\subsection{驱动因素5：[名称]}
[详细分析]

\section{PESTLE分析}

% 视觉：PESTLE分析图
% python skills/scientific-schematics/scripts/generate_schematic.py "PESTLE分析图，中心六边形标注市场，周围6个六边形：政治(红)、经济(蓝)、社会(绿)、技术(橙)、法律(紫)、环境(蓝绿)。每个外围含2-3个关键因素要点" -o figures/pestle_analysis.png --doc-type report

\begin{figure}[htbp]
\centering
% \includegraphics[width=0.9\textwidth]{../figures/pestle_analysis.png}
\caption{PESTLE分析框架}
\label{fig:pestle}
\end{figure}

\subsection{政治因素}
[分析影响市场的政治因素：
- 政府政策
- 政治稳定性
- 贸易政策
- 税收法规]

\subsection{经济因素}
[分析经济因素：
- 经济增长
- 利率
- 通货膨胀
- 汇率
- 消费者支出]

\subsection{社会因素}
[分析社会因素：
- 人口统计
- 文化趋势
- 消费者态度
- 劳动力趋势]

\subsection{技术因素}
[分析技术因素：
- 技术采用
- 研发活动
- 自动化
- 数字化转型]

\subsection{法律因素}
[分析法律因素：
- 行业法规
- 合规要求
- 知识产权
- 劳动法]

\subsection{环境因素}
[分析环境因素：
- 可持续性要求
- 环境法规
- 气候影响
- 资源可用性]

\section{新兴趋势}

% 视觉：趋势时间线
% python skills/scientific-schematics/scripts/generate_schematic.py "水平时间线展示2024至2030年新兴趋势。在不同时间点标注6-8个趋势含图标和标签。技术趋势蓝色，市场趋势绿色，监管趋势橙色。每个趋势下方简要描述" -o figures/emerging_trends_timeline.png --doc-type report

\begin{figure}[htbp]
\centering
% \includegraphics[width=\textwidth]{../figures/trends_timeline.png}
\caption{新兴行业趋势时间线}
\label{fig:trends}
\end{figure}

[识别和分析将塑造市场的新兴趋势：]

\subsection{趋势1：[名称]}
[详细趋势分析]

\subsection{趋势2：[名称]}
[详细趋势分析]

\subsection{趋势3：[名称]}
[详细趋势分析]

\section{增长抑制因素}

[识别可能减缓市场增长的因素：
- 市场壁垒
- 资源限制
- 采用挑战
- 竞争压力]

% ============================================================================
% 第4章：竞争格局（6-8页）
% ============================================================================
\chapter{竞争格局}

\section{市场结构分析}

[分析市场的竞争结构：]

\subsection{市场集中度}

[提供市场集中度分析：
- 竞争者数量
- 市场份额分布
- 集中度指标（HHI、CR4）]

\section{波特五力分析}

% 视觉：波特五力
% python skills/scientific-schematics/scripts/generate_schematic.py "波特五力分析图。中心框标注同业竞争[高/中/低]。四周框箭头连接：新进入者威胁[评级](上)、供应商议价能力[评级](左)、买方议价能力[评级](右)、替代品威胁[评级](下)。颜色编码高=红、中=橙、低=绿" -o figures/porters_five_forces.png --doc-type report

\begin{figure}[htbp]
\centering
% \includegraphics[width=0.85\textwidth]{../figures/porters_five_forces.png}
\caption{波特五力分析}
\label{fig:porter}
\end{figure}

\begin{porterbox}[波特五力总结]
\begin{tabular}{@{}ll@{}}
\textbf{力量} & \textbf{评级} \\
\midrule
新进入者威胁 & \riskmedium{} \\
供应商议价能力 & \risklow{} \\
买方议价能力 & \riskhigh{} \\
替代品威胁 & \risklow{} \\
竞争对抗 & \riskhigh{} \\
\end{tabular}
\end{porterbox}

\subsection{新进入者威胁}
[对进入壁垒和威胁水平的详细分析]

\subsection{供应商议价能力}
[供应商力量动态的详细分析]

\subsection{买方议价能力}
[买方力量动态的详细分析]

\subsection{替代品威胁}
[替代产品/服务的详细分析]

\subsection{竞争对抗}
[竞争强度的详细分析]

\section{市场份额分析}

% 视觉：市场份额图
% python skills/scientific-schematics/scripts/generate_schematic.py "饼图展示前10家企业市场份额。企业A XX%，企业B XX%，企业C XX%，其他XX%。各企业不同颜色。含企业名称和百分比图例" -o figures/market_share.png --doc-type report

\begin{figure}[htbp]
\centering
% \includegraphics[width=0.8\textwidth]{../figures/market_share.png}
\caption{各公司市场份额（2024年）}
\label{fig:market_share}
\end{figure}

\begin{table}[htbp]
\centering
\caption{按市场份额排名前十的公司}
\begin{tabular}{@{}clrrr@{}}
\toprule
\textbf{排名} & \textbf{公司} & \textbf{营收} & \textbf{份额} & \textbf{同比增长} \\
\midrule
1 & 公司A & \$X.XX十亿 & XX.X\% & \trendup{} XX\% \\
\rowcolor{tablealt} 2 & 公司B & \$X.XX十亿 & XX.X\% & \trendup{} XX\% \\
3 & 公司C & \$X.XX十亿 & XX.X\% & \trendflat{} XX\% \\
\rowcolor{tablealt} 4 & 公司D & \$X.XX十亿 & XX.X\% & \trendup{} XX\% \\
5 & 公司E & \$X.XX十亿 & XX.X\% & \trenddown{} XX\% \\
\rowcolor{tablealt} 6 & 公司F & \$X.XX十亿 & XX.X\% & \trendup{} XX\% \\
7 & 公司G & \$X.XX十亿 & XX.X\% & \trendup{} XX\% \\
\rowcolor{tablealt} 8 & 公司H & \$X.XX十亿 & XX.X\% & \trendflat{} XX\% \\
9 & 公司I & \$X.XX十亿 & XX.X\% & \trendup{} XX\% \\
\rowcolor{tablealt} 10 & 公司J & \$X.XX十亿 & XX.X\% & \trenddown{} XX\% \\
\midrule
& 其他 & \$X.XX十亿 & XX.X\% & -- \\
\bottomrule
\end{tabular}
\label{tab:market_share}
\end{table}

\section{竞争定位}

% 视觉：竞争定位矩阵
% python skills/scientific-schematics/scripts/generate_schematic.py "2x2竞争定位矩阵。X轴市场聚焦利基到广泛，Y轴解决方案产品到平台。绘制8-10个企业为不同大小标注圆圈(代表市场份额)。含象限标签" -o figures/competitive_positioning.png --doc-type report

\begin{figure}[htbp]
\centering
% \includegraphics[width=0.9\textwidth]{../figures/competitive_positioning.png}
\caption{竞争定位矩阵}
\label{fig:competitive_positioning}
\end{figure}

[分析竞争者在关键维度上的市场定位：]

\subsection{战略集团}

% 视觉：战略集团图
% python skills/scientific-schematics/scripts/generate_schematic.py "战略群组图，圆圈代表不同战略群组。X轴地理范围区域到全球，Y轴产品广度窄到宽。每个圆圈含多个企业名称，大小为集体市场份额。4-5个群组" -o figures/strategic_group_map.png --doc-type report

\begin{figure}[htbp]
\centering
% \includegraphics[width=0.85\textwidth]{../figures/strategic_groups.png}
\caption{战略集团映射}
\label{fig:strategic_groups}
\end{figure}

[识别和描述竞争格局中的战略集团]

\section{竞争动态}

[分析竞争行为和动态：
- 近期并购活动
- 合作伙伴关系公告
- 产品发布
- 定价趋势
- 地理扩张]

\section{进入壁垒}

[分析保护在位者并挑战新进入者的壁垒：
- 资本要求
- 监管壁垒
- 技术壁垒
- 品牌和声誉
- 分销渠道获取
- 规模经济]

% ============================================================================
% 第5章：客户分析与细分（4-5页）
% ============================================================================
\chapter{客户分析 \& 细分}

\section{客户细分}

% 视觉：客户细分
% python skills/scientific-schematics/scripts/generate_schematic.py "树状图或饼图展示客户细分。细分A XX%(大型企业)，细分B XX%(中型企业)，细分C XX%(小型企业)，细分D XX%(其他)。大小代表市场份额。不同颜色" -o figures/customer_segments.png --doc-type report

\begin{figure}[htbp]
\centering
% \includegraphics[width=0.85\textwidth]{../figures/customer_segments.png}
\caption{按市场份额划分的客户细分}
\label{fig:customer_segments}
\end{figure}

\begin{table}[htbp]
\centering
\caption{客户细分分析}
\begin{tabular}{@{}lrrrr@{}}
\toprule
\textbf{细分市场} & \textbf{规模} & \textbf{增长率} & \textbf{平均交易额} & \textbf{客户获取成本} \\
\midrule
大型企业 & \$X.XX十亿 & XX\% & \$XXX千 & \$XX千 \\
\rowcolor{tablealt} 中型市场 & \$X.XX十亿 & XX\% & \$XX千 & \$X千 \\
中小企业 & \$X.XX十亿 & XX\% & \$X千 & \$X千 \\
\rowcolor{tablealt} 消费者 & \$X.XX十亿 & XX\% & \$XXX & \$XX \\
\bottomrule
\end{tabular}
\label{tab:customer_segments}
\end{table}

\subsection{细分市场A：[大型企业]}
[详细细分市场分析，包括：
- 细分市场特征
- 购买行为
- 关键需求和痛点
- 决策过程
- 支付意愿]

\subsection{细分市场B：[中型市场]}
[详细细分市场分析]

\subsection{细分市场C：[中小企业]}
[详细细分市场分析]

\subsection{细分市场D：[其他]}
[详细细分市场分析]

\section{细分市场吸引力}

% 视觉：细分市场吸引力矩阵
% python skills/scientific-schematics/scripts/generate_schematic.py "2x2细分吸引力矩阵。X轴细分规模小到大，Y轴增长率低到高。绘制客户细分为圆圈。右上'优先重点'，左上'投资增长'，右下'收获'，左下'降低优先'。含细分标签" -o figures/segment_attractiveness.png --doc-type report

\begin{figure}[htbp]
\centering
% \includegraphics[width=0.85\textwidth]{../figures/segment_attractiveness.png}
\caption{客户细分吸引力矩阵}
\label{fig:segment_attractiveness}
\end{figure}

[分析哪些细分市场对投资和聚焦最具吸引力]

\section{客户需求分析}

[识别和按细分市场优先排序客户需求：
- 功能需求
- 情感需求
- 社交需求
- 痛点
- 未满足需求]

\section{购买行为}

[分析该市场的客户购买方式：
- 购买触发因素
- 决策过程
- 关键影响者
- 评估标准
- 购买渠道]

% 视觉：客户旅程
% python skills/scientific-schematics/scripts/generate_schematic.py "客户旅程图展示5个阶段：认知→考虑→决策→实施→拥护。每个阶段展示关键活动、痛点和触点。水平流程每阶段配图标" -o figures/customer_journey.png --doc-type report

\begin{figure}[htbp]
\centering
% \includegraphics[width=\textwidth]{../figures/customer_journey.png}
\caption{客户旅程图}
\label{fig:customer_journey}
\end{figure}

% ============================================================================
% 第6章：技术与创新格局（4-5页）
% ============================================================================
\chapter{技术 \& 创新格局}

\section{当前技术栈}

[描述市场中常用的技术基础设施和栈]

\section{技术路线图}

% 视觉：技术路线图
% python skills/scientific-schematics/scripts/generate_schematic.py "2024至2030技术路线图。3条平行轨道：核心技术(蓝)、新兴技术(绿)、使能技术(橙)。各轨道标注关键里程碑和技术引入节点。水平时间线格式" -o figures/technology_roadmap.png --doc-type report

\begin{figure}[htbp]
\centering
% \includegraphics[width=\textwidth]{../figures/technology_roadmap.png}
\caption{技术演进路线图（2024-2030）}
\label{fig:tech_roadmap}
\end{figure}

[描述技术预期如何演进：
- 近期（1-2年）
- 中期（3-5年）
- 远期（5-10年）]

\section{新兴技术}

[分析可能影响市场的新兴技术：
- 技术描述
- 当前成熟度
- 预计主流采用时间线
- 对市场的潜在影响]

\subsection{技术1：[名称]}
[详细分析]

\subsection{技术2：[名称]}
[详细分析]

\subsection{技术3：[名称]}
[详细分析]

\section{创新趋势}

% 视觉：创新矩阵
% python skills/scientific-schematics/scripts/generate_schematic.py "创新采用曲线或炒作周期图展示关键技术所处阶段。从左到右：创新萌芽、过高期望峰值、幻灭低谷、稳步爬升、生产成熟期。在不同位置绘制6-8个技术" -o figures/innovation_adoption_curve.png --doc-type report

\begin{figure}[htbp]
\centering
% \includegraphics[width=\textwidth]{../figures/innovation_curve.png}
\caption{技术采用曲线/炒作周期}
\label{fig:innovation}
\end{figure}

[分析创新趋势和研发活动：
- 研发投入水平
- 专利申请趋势
- 创业活动
- 企业创新举措]

\section{技术采用壁垒}

[识别技术采用的壁垒：
- 技术复杂性
- 集成挑战
- 成本壁垒
- 技能差距
- 安全/隐私顾虑]

% ============================================================================
% 第7章：监管与政策环境（3-4页）
% ============================================================================
\chapter{监管 \& 政策环境}

\section{当前监管框架}

[描述当前监管格局：
- 关键法规
- 监管机构
- 合规要求
- 执法机制]

\section{监管时间线}

% 视觉：监管时间线
% python skills/scientific-schematics/scripts/generate_schematic.py "2020至2028监管时间线。水平时间线上标注关键监管里程碑。过去事件深蓝色，未来/即将出台浅蓝色。含法规名称和生效日期。当前日期垂直线标记" -o figures/regulatory_timeline.png --doc-type report

\begin{figure}[htbp]
\centering
% \includegraphics[width=\textwidth]{../figures/regulatory_timeline.png}
\caption{监管发展时间线}
\label{fig:regulatory_timeline}
\end{figure}

[按时间顺序分析监管发展]

\section{监管影响分析}

[分析法规如何影响市场：
- 合规成本
- 市场准入影响
- 竞争影响
- 产品/服务要求]

\section{政策趋势}

[识别可能影响市场的政策趋势：
- 政府优先事项
- 资金计划
- 贸易政策
- 环境政策]

\section{区域监管差异}

[比较各区域的监管环境：
- 北美
- 欧洲
- 亚太
- 其他地区]

% ============================================================================
% 第8章：风险分析（3-4页）
% ============================================================================
\chapter{风险分析}

\section{风险概述}

[概述市场参与者面临的关键风险]

\section{风险评估}

% 视觉：风险热力图
% python skills/scientific-schematics/scripts/generate_schematic.py "风险热力图矩阵。X轴影响低到严重，Y轴概率不太可能到非常可能。绘制10-12个风险标注圆点。颜色编码：绿色低风险、黄色中等、橙色高风险、红色严重。含风险标签" -o figures/risk_heatmap.png --doc-type report

\begin{figure}[htbp]
\centering
% \includegraphics[width=0.9\textwidth]{../figures/risk_heatmap.png}
\caption{风险评估热力图}
\label{fig:risk_heatmap}
\end{figure}

\begin{table}[htbp]
\centering
\caption{风险登记摘要}
\begin{tabular}{@{}llccl@{}}
\toprule
\textbf{风险} & \textbf{类别} & \textbf{概率} & \textbf{影响} & \textbf{评级} \\
\midrule
风险1 & 市场 & 高 & 高 & \riskhigh{} \\
\rowcolor{tablealt} 风险2 & 监管 & 中 & 高 & \riskhigh{} \\
风险3 & 技术 & 中 & 中 & \riskmedium{} \\
\rowcolor{tablealt} 风险4 & 竞争 & 高 & 中 & \riskmedium{} \\
风险5 & 运营 & 低 & 高 & \riskmedium{} \\
\rowcolor{tablealt} 风险6 & 财务 & 低 & 中 & \risklow{} \\
\bottomrule
\end{tabular}
\label{tab:risk_register}
\end{table}

\subsection{市场风险}
[市场相关风险的详细分析]

\subsection{竞争风险}
[竞争风险的详细分析]

\subsection{监管风险}
[监管风险的详细分析]

\subsection{技术风险}
[技术风险的详细分析]

\subsection{运营风险}
[运营风险的详细分析]

\subsection{财务风险}
[财务风险的详细分析]

\section{风险缓解策略}

% 视觉：风险缓解矩阵
% python skills/scientific-schematics/scripts/generate_schematic.py "风险应对矩阵，左列为风险，右列为对应应对策略。箭头连接风险与应对。按风险严重程度着色。含预防和响应策略" -o figures/risk_mitigation.png --doc-type report

\begin{figure}[htbp]
\centering
% \includegraphics[width=0.9\textwidth]{../figures/risk_mitigation.png}
\caption{风险缓解框架}
\label{fig:risk_mitigation}
\end{figure}

[描述已识别风险的缓解策略]

\begin{riskbox}[风险缓解总结]
\begin{enumerate}
    \item \textbf{[风险1：]} [缓解策略]
    \item \textbf{[风险2：]} [缓解策略]
    \item \textbf{[风险3：]} [缓解策略]
    \item \textbf{[风险4：]} [缓解策略]
\end{enumerate}
\end{riskbox}

% ============================================================================
% 第9章：战略机会与建议（4-5页）
% ============================================================================
\chapter{战略机会 \& 建议}

\section{机会分析}

% 视觉：机会矩阵
% python skills/scientific-schematics/scripts/generate_schematic.py "2x2机会矩阵。X轴市场吸引力低到高，Y轴胜出能力低到高。绘制6-8个机会为不同大小标注圆圈。右上'积极进取'，左上'选择性投资'，右下'能力建设'，左下'回避'" -o figures/opportunity_matrix.png --doc-type report

\begin{figure}[htbp]
\centering
% \includegraphics[width=0.9\textwidth]{../figures/opportunity_matrix.png}
\caption{战略机会评估矩阵}
\label{fig:opportunity_matrix}
\end{figure}

[识别和分析市场中的战略机会]

\subsection{机会1：[名称]}
\begin{opportunitybox}[机会：[名称]]
\textbf{描述：} [简要描述]

\textbf{市场规模：} \marketsize{X.X十亿美元}

\textbf{增长率：} \growthrate{XX.X}

\textbf{战略匹配度：} \rating{4}

\textbf{所需投资：} \$XX百万
\end{opportunitybox}

[机会的详细分析]

\subsection{机会2：[名称]}
[详细分析]

\subsection{机会3：[名称]}
[详细分析]

\section{战略选项分析}

[分析不同的战略路径：
- 自建（有机发展）
- 收购（并购）
- 合作（战略联盟）
- 放弃（不追求）]

\section{优先排序建议}

% 视觉：建议优先级矩阵
% python skills/scientific-schematics/scripts/generate_schematic.py "优先级矩阵展示建议。X轴投入低到高，Y轴影响低到高。绘制6-8个建议。左上'速赢项'，右上'重大项目'，左下'填充项'，右下'吃力不讨好'。标注每个建议" -o figures/priority_matrix.png --doc-type report

\begin{figure}[htbp]
\centering
% \includegraphics[width=0.85\textwidth]{../figures/recommendation_priority.png}
\caption{建议优先级框架}
\label{fig:recommendations}
\end{figure}

\begin{recommendationbox}[战略建议]
\textbf{第一优先级：立即行动}
\begin{enumerate}
    \item \textbf{[建议1：]} [详细行动及预期成果、时间线和投资]
    \item \textbf{[建议2：]} [详细行动]
\end{enumerate}

\textbf{第二优先级：近期（6-12个月）}
\begin{enumerate}[start=3]
    \item \textbf{[建议3：]} [详细行动]
    \item \textbf{[建议4：]} [详细行动]
\end{enumerate}

\textbf{第三优先级：中期（1-2年）}
\begin{enumerate}[start=5]
    \item \textbf{[建议5：]} [详细行动]
    \item \textbf{[建议6：]} [详细行动]
\end{enumerate}
\end{recommendationbox}

\section{成功因素}

[识别实施建议的关键成功因素：
- 组织能力
- 资源要求
- 时间考量
- 外部依赖]

% ============================================================================
% 第10章：实施路线图（3-4页）
% ============================================================================
\chapter{实施路线图}

\section{实施概述}

[概述实施方法和时间线]

\section{分阶段实施计划}

% 视觉：实施时间线
% python skills/scientific-schematics/scripts/generate_schematic.py "甘特图风格实施时间线展示24个月4个阶段。阶段1基础建设(1-6月)、阶段2构建开发(4-12月)、阶段3规模扩展(10-18月)、阶段4优化完善(16-24月)。各阶段重叠，关键里程碑标记。各阶段不同颜色" -o figures/implementation_timeline.png --doc-type report

\begin{figure}[htbp]
\centering
% \includegraphics[width=\textwidth]{../figures/implementation_timeline.png}
\caption{实施路线图时间线}
\label{fig:implementation}
\end{figure}

\subsection{第一阶段：基础建设（第1-6个月）}
[第一阶段的详细活动和交付物]

\subsection{第二阶段：构建（第4-12个月）}
[第二阶段的详细活动和交付物]

\subsection{第三阶段：扩展（第10-18个月）}
[第三阶段的详细活动和交付物]

\subsection{第四阶段：优化（第16-24个月）}
[第四阶段的详细活动和交付物]

\section{关键里程碑}

% 视觉：里程碑跟踪器
% python skills/scientific-schematics/scripts/generate_schematic.py "里程碑追踪展示水平时间线上8-10个关键里程碑。每个里程碑含日期、名称和状态指示(已完成=绿色勾选、进行中=黄色圆圈、待开始=灰色圆圈)。按阶段分组" -o figures/milestone_tracker.png --doc-type report

\begin{figure}[htbp]
\centering
% \includegraphics[width=\textwidth]{../figures/milestone_tracker.png}
\caption{关键实施里程碑}
\label{fig:milestones}
\end{figure}

\begin{table}[htbp]
\centering
\caption{实施里程碑}
\begin{tabular}{@{}llll@{}}
\toprule
\textbf{里程碑} & \textbf{目标日期} & \textbf{负责人} & \textbf{成功标准} \\
\midrule
里程碑1 & 第3个月 & [负责人] & [标准] \\
\rowcolor{tablealt} 里程碑2 & 第6个月 & [负责人] & [标准] \\
里程碑3 & 第9个月 & [负责人] & [标准] \\
\rowcolor{tablealt} 里程碑4 & 第12个月 & [负责人] & [标准] \\
里程碑5 & 第18个月 & [负责人] & [标准] \\
\rowcolor{tablealt} 里程碑6 & 第24个月 & [负责人] & [标准] \\
\bottomrule
\end{tabular}
\label{tab:milestones}
\end{table}

\section{资源需求}

[详细说明实施的资源需求：
- 团队结构
- 预算分配
- 技术需求
- 外部支持需求]

\section{治理结构}

[定义实施治理：
- 决策权限
- 汇报结构
- 审查节奏
- 升级路径]

% ============================================================================
% 第11章：投资论证与财务预测（3-4页）
% ============================================================================
\chapter{投资论证 \& 财务预测}

\section{投资概述}

[总结投资机会：
- 关键价值驱动因素
- 预期回报
- 投资时间线
- 风险调整评估]

\begin{executivesummarybox}[投资论证]
\reporttitle{} 市场呈现出引人注目的投资机会，具有以下特点：

\begin{itemize}
    \item \textbf{庞大市场：} \marketsize{XX十亿美元} TAM以 \growthrate{XX.X} CAGR增长
    \item \textbf{有利动态：} [关键市场动态]
    \item \textbf{强劲驱动：} [关键增长驱动因素]
    \item \textbf{可控风险：} [风险概述]
    \item \textbf{清晰价值路径：} [价值创造概述]
\end{itemize}
\end{executivesummarybox}

\section{财务预测}

% 视觉：财务预测
% python skills/scientific-schematics/scripts/generate_schematic.py "财务预测图展示5年收入增长。柱状图为收入叠折线为增长率。三种情景：保守(灰)、基准(蓝)、乐观(绿)。双Y轴：收入($M)和增长率(%)。含数据标签" -o figures/financial_projections.png --doc-type report

\begin{figure}[htbp]
\centering
% \includegraphics[width=0.9\textwidth]{../figures/financial_projections.png}
\caption{财务预测（5年）}
\label{fig:financials}
\end{figure}

\begin{table}[htbp]
\centering
\caption{财务预测摘要}
\begin{tabular}{@{}lrrrrr@{}}
\toprule
\textbf{指标} & \textbf{第1年} & \textbf{第2年} & \textbf{第3年} & \textbf{第4年} & \textbf{第5年} \\
\midrule
营收（\$M） & \$XX & \$XX & \$XX & \$XX & \$XX \\
\rowcolor{tablealt} 增长率 & XX\% & XX\% & XX\% & XX\% & XX\% \\
毛利率 & XX\% & XX\% & XX\% & XX\% & XX\% \\
\rowcolor{tablealt} EBITDA（\$M） & \$XX & \$XX & \$XX & \$XX & \$XX \\
EBITDA利润率 & XX\% & XX\% & XX\% & XX\% & XX\% \\
\bottomrule
\end{tabular}
\label{tab:financials}
\end{table}

\section{情景分析}

% 视觉：情景比较
% python skills/scientific-schematics/scripts/generate_schematic.py "情景对比图展示3种情景(保守、基准、乐观)各关键指标。分组柱状图，X轴指标(第5年收入、第5年EBITDA、市场份额、ROI)，Y轴数值。按情景着色" -o figures/scenario_comparison.png --doc-type report

\begin{figure}[htbp]
\centering
% \includegraphics[width=0.85\textwidth]{../figures/scenario_analysis.png}
\caption{情景分析比较}
\label{fig:scenarios}
\end{figure}

\subsection{保守情景}
[保守情景的详细假设和结果]

\subsection{基准情景}
[基准情景的详细假设和结果]

\subsection{乐观情景}
[乐观情景的详细假设和结果]

\section{关键假设}

[记录财务预测所依据的关键假设：
- 市场增长假设
- 定价假设
- 成本假设
- 竞争假设
- 时间假设]

\section{敏感性分析}

[分析预测对关键变量的敏感性：
- 营收对市场增长的敏感性
- 利润率对定价的敏感性
- 回报对时间的敏感性]

\section{回报预期}

[总结预期回报：
- ROI预测
- 回收期
- IRR估计
- 倍数分析]

% ============================================================================
% 附录
% ============================================================================
\appendix

% ============================================================================
% 附录A：研究方法与数据来源
% ============================================================================
\chapter{研究方法 \& 数据来源}

\section{研究方法论}

[描述所使用的研究方法：
- 一手研究方法
- 二手研究来源
- 数据收集时间框架
- 应用的分析框架]

\section{数据来源}

[列出报告中使用的所有数据来源：
- 市场研究报告
- 行业数据库
- 政府统计数据
- 企业报告
- 专家访谈
- 学术出版物]

\section{局限性}

[承认分析的局限性：
- 数据可用性限制
- 方法论局限
- 预测不确定性
- 范围限制]

% ============================================================================
% 附录B：详细市场数据
% ============================================================================
\chapter{详细市场数据}

\section{历史市场数据}

[提供详细的历史市场数据表]

\section{区域数据细分}

[提供详细的区域市场数据]

\section{细分市场数据详情}

[提供详细的细分市场数据]

\section{竞争数据}

[提供详细的竞争数据表]

% ============================================================================
% 附录C：企业概况
% ============================================================================
\chapter{企业概况}

\section{公司A}

\begin{calloutbox}[企业概况：公司A]
\textbf{总部：} [地点]

\textbf{营收：} \$X.X十亿美元（2024财年）

\textbf{员工数：} X,XXX

\textbf{市场地位：} [地位描述]

\textbf{主要产品/服务：} [列表]

\textbf{近期发展：} [概述]
\end{calloutbox}

[公司战略和定位的简要叙述]

\section{公司B}
[企业概况]

\section{公司C}
[企业概况]

\section{公司D}
[企业概况]

\section{公司E}
[企业概况]

% ============================================================================
% 参考文献
% ============================================================================
\newpage
\bibliographystyle{plainnat}
\bibliography{../references/references}

% 替代方案：不使用BibTeX时手动添加参考文献
% \begin{thebibliography}{99}
% 
% \bibitem{source1}
% Author1, A.B. (2024). 
% Title of report or article. 
% \textit{Publisher/Source}. 
% URL
% 
% \bibitem{source2}
% [继续列出所有参考文献...]
%
% \end{thebibliography}

% ============================================================================
% 文档结束
% ============================================================================
\end{document}